\section{Exact scaling under a controlled row defect}
\label{sec:scaling}
For a vector \(x\), write \(\|x\|_\infty=\max_i|x_i|\);
matrix operator norms are Euclidean unless specified.
Let \(s=\one_A-\one_B\).

\begin{lemma}[A quantitative bipartite scaling repair]
\label{lem:scaling}
Fix \(\delta\in(0,1)\).
Let \(Q\) be a symmetric nonnegative matrix supported between
two parts of equal size, and put
\[
 d=Q\one,\qquad D=\diag(d),\qquad
 \|d-\one\|_\infty\leq\varepsilon<1/2,\qquad
 S=D^{-1/2}QD^{-1/2}.
\]
Assume all eigenvalues of \(S\) except its simple eigenvalues
\(1,-1\) have modulus at most \(1-\delta\).
There are constants \(c_\delta,C_\delta>0\) such that if
\(\varepsilon L^2\leq c_\delta\), there is a positive vector
\(u\) with
\[
 P=\diag(u)Q\diag(u),\quad P\one=\one,\qquad
 \|u-\one\|_\infty\leq C_\delta\varepsilon L,\qquad
 \|P-Q\|\leq C_\delta\varepsilon L.                           \tag{5}
\]
The support of \(P\) is exactly the support of \(Q\).
\end{lemma}
\begin{proof}
The bipartite symmetry gives
\(\sum_A d_i=\sum_Bd_i\).
Set \(T=D^{-1}Q\), a reversible stochastic matrix with stationary
measure \(\pi_i=d_i/\sum_jd_j\).
It is similar to \(S\), and
\(\pi(A)=\pi(B)=1/2\).
Its eigenvectors at \(1,-1\) are \(\one,s\).
The corresponding projections are
\[
 \Pi_+f=\one\sum_i\pi_i f_i,\qquad
 \Pi_-f=s\sum_i\pi_i s_i f_i,\qquad
 \Pi_0=I-\Pi_+-\Pi_-.
\]
For \(j\geq0\), the assumptions and \(d_i\in[1/2,3/2]\) give
\[
 \|T^j\Pi_0\|_{\infty\to\infty}
 \leq C\min\{1,\sqrt N(1-\delta)^j\}.                         \tag{6}
\]
For the first bound, use \(T^j\Pi_0=T^j-\Pi_+-(-1)^j\Pi_-\)
and the three bounded stochastic/projection norms.
For the second, use the \(L^2(\pi)\) spectral gap and
\(\min_i\pi_i\geq1/(3N)\).

Let \({\cal L}_Q=D+Q=D(I+T)\).
Its kernel is \(\R s\), by the hypothesis at \(-1\).
For \(r\perp s\), the vector \(f=D^{-1}r\) has
\(\Pi_-f=0\).
The convergent series
\[
 z_0=\tfrac12\Pi_+f+
                 \sum_{j\geq0}(-1)^jT^j\Pi_0f
\]
therefore solves \({\cal L}_Qz_0=r\).
The sum of the bounds in (6) is \(O_\delta(L)\).
Subtract the \(s\)-component of \(z_0\) to impose \(z\perp s\).
Thus the inverse on \(s^\perp\) satisfies
\[
 \|{\cal L}_Q^\dagger r\|_\infty
                    \leq C_\delta L\|r\|_\infty.             \tag{7}
\]
Here the dagger means the inverse with this fixed gauge;
no inverse is asserted on the one-dimensional kernel.

Put \(u=\one+z\).
The exact row equation, including its nonlinear cross term, is
\[
 u\odot(Qu)-\one
 =(d-\one)+{\cal L}_Qz+z\odot(Qz).                           \tag{8}
\]
Both \(d-\one\) and \(z\odot(Qz)\) are orthogonal to \(s\).
For the latter,
\[
 \sum_{i\in A}z_i(Qz)_i
 =\sum_{i\in A,j\in B}Q_{ij}z_i z_j
 =\sum_{j\in B}z_j(Qz)_j.
\]
For the former use equal part sizes and equal total degrees.
Consequently the map
\[
 z\longmapsto-{\cal L}_Q^\dagger
                    \bigl((d-\one)+z\odot(Qz)\bigr)
\]
maps \(s^\perp\) to itself.
By (7), it is a contraction on the ball
\(\|z\|_\infty\leq2C_\delta\varepsilon L\) when
\(c_\delta\) is sufficiently small.
Indeed \(\|Q\|_{\infty\to\infty}\leq3/2\), and the nonlinear
Lipschitz constant on a ball of radius \(r\) is at most \(3r\);
the contraction cost is \(O_\delta(\varepsilon L^2)\).
The map also preserves this ball by the same inequality.
The fixed point gives (8) equal to zero.
Choose \(c_\delta\) smaller so that \(\|z\|_\infty<1/2\);
then every \(u_i\) is positive.

Finally, with \(Z=\diag(z)\), retain all terms
\[
 P-Q=ZQ+QZ+ZQZ.
\]
Since \(\|Q\|\leq3/2\), their operator norm is
\(O_\delta(\varepsilon L)\), proving (5).
\end{proof}

\begin{corollary}[Repair after forbidden edges]
\label{cor:forbiddenrepair}
Let \(P\) satisfy the hypotheses of Theorem~\ref{thm:rounding},
and let \(F\) be an arbitrary set of its support edges.
Define the actual removed row mass
\[
 \eta=\max_i\sum_{j:ij\in F}P_{ij}.
\]
If \(\eta L^2\leq c_\delta\), then there is a symmetric
stochastic \(P_F\) supported exactly on the unused edges, with
\[
 \|P_F-P\|\leq C_\delta\eta L,\qquad
 \max_{ij}(P_F)_{ij}\leq2\alpha.                             \tag{9}
\]
It has nontrivial gap at least \(\delta/2\).
\end{corollary}
\begin{proof}
Set \(Q_{ij}=P_{ij}\one_{ij\notin F}\).
Then \(Q\one\) lies between \(1-\eta\) and 1.
Symmetry and the row bound give
\(\|Q-P\|\leq\eta\).
For \(D=\diag(Q\one)\), it follows that
\[
 \|D^{-1/2}QD^{-1/2}-P\|\leq C\eta.
\]
The normalized matrix has eigenvalues \(1,-1\), and Weyl's
inequality bounds every other eigenvalue in modulus by
\(1-\delta+C\eta\).
Thus Lemma~\ref{lem:scaling} applies with gap \(\delta/2\).
Its exact positive scaling preserves all unused entries and
has multipliers \(1+O_\delta(\eta L)\).
This gives (9).
After scaling, the constant and parity vectors are again exact
eigenvectors. The operator perturbation in (9), and a further
decrease of \(c_\delta\), give the claimed gap.
\end{proof}
